\chapter{绪论}

\section{研究背景}
示例引用公式~\ref{gt1}和公式~\ref{gt2}
\begin{equation}
 C^f_t=\left \{a(P_t^{gen})^2 + b P_t^{gen}+c\right \} \bigtriangleup t
 \label{gt1}
\end{equation}
公式原则上应居中书写。若公式前有文字（如“解”、“假定”等），文字空两格写，公式仍居中写。公式末不加标点。
公式序号按章编排，如第1章第一个公式序号为“(1-1)”。
文中引用公式时，一般用“见式 (1-1)”或“由公式 (1-1)”。
公式中用斜线表示“除”的关系时应采用括号，以免含糊不清。通常“乘”的关系在前。不能用文字形式表示等式。
\begin{gather}
    C^f_t=\left \{a(P_t^{gen})^2 + b P_t^{gen}+c\right \} \bigtriangleup t \label{gt2} \\
    C^f_t=\left \{a(P_t^{gen})^2 + b P_t^{gen}+c\right \} \bigtriangleup t \label{gt3}
\end{gather}
式中~~  $P_t^{gen}$------ 发电量；\breakwithtilde $P_t^{gen}$------ 发电量；
\breakwithtilde $P_t^{gen}$------ 发电量。


